\head{第二十三版整合核查：第15节的实际范围}
\audit{V23-INT01}{连续源页与实际接续}{已有公理化草稿的法英回校，非重复首译}
本轮从已发布第二十二版的实际源码接续，基线提交为\texttt{a1e4d5b5}。上一轮法英回校确实止于第一篇第14节；第15节已有完整中文译文，保存在\path{continuations/v22_section15_inherited.tex}，共14,614字节。本版保留这一原文件，按法文手稿及英文编订稿重新核对其全部段落，从第15节标题直至给收信人的最后建议。

法文连续范围为完整扫描的物理90页末至94页末；英文范围为物理59页节标题至61页末。正文分为公理化设置与光滑性条件两模块，按原顺序接合；源页里的研究建议、推测、写作指示、公式与交换图各按其身份保留。法稿第95页和英稿第62页已进入第二篇第16节，不计入本轮新回校范围。

对应的两个专题报告逐项说明法文原句、英文转写、旧中文判断及本版处理。法文第91页关于正式成稿最初几节的写作指示、第92--93页的像集与非空性、从(2.1)至(2.11)的范围表述，以及第94页关于(2.12)至(2.18)的句子，都须按实际原页阅读，不能把英文省略或误排归为原作者的数学错误。

\audit{V23-INT02}{正文订正与原稿意见的边界}{公式、论证对象和适用条件分别核查}
\lead{泛截面的维数必须逐不可约分支计算。}
对于可能混维的闭集，仅有最大的$\dim Z$及$\dim f(Z)\ge m$，不足以推出原先笼统的$\dim Z-m$。高维分支可以被压到维数小于$m$的像而被泛截面避开，留下的低维分支才决定实际截面维数。本版正文采用对像维数至少$m$的分支取最大值的表达，专题报告给出高维压塌分支与嵌入曲线的反例，以及在明确有限型条件下的证明。非空性公理本身与维数结论仍分开陈述，不把公理写成仅由光滑性自动推出的定理。

\lead{初始公理不足以保持所有局部等价。}
初段并未要求$P$光滑，而几何不可约也不排除$G$带有幂零元。取$P=G=\mathcal P=X=\operatorname{Spec}k[\epsilon]/(\epsilon^2)$，投影$q,p$及$f$为恒等，$N=m=0$，初始公理仍成立；但泛截面映射$\mathcal X_\eta=\operatorname{Spec}k\to X$不平坦，泛截面正则而$X$不正则。因此不能由初始公理直接宣布(2.1)至(2.11)的全部局部等价照搬。正文把依赖$h\colon\mathcal X_\eta\to X$正则性的推广明确列为需补条件，报告给出$G$在泛点处约化这一充分版本。后文针对光滑$P$、$X$的泛截面光滑性修复是另一论证，不证明任意非光滑$X$的全部历史命题。

\lead{非分歧条件须用正确的态射。}
本节中关联预概形到环境$P$的投影，与其到$P\times G$的态射及逐参数纤维态射不是同一论证对象。用于截面切空间单射的条件针对后者；把前一个具有正相对维数的光滑投影当作应当非分歧的态射，会错误地扩大坏点集。原稿写法、相应订正及证明中的用途分别列于脚注和专题报告。

\lead{光滑坏集的维数估计需处理非约化参数空间。}
原假设允许$G$非约化，因而不能未经平坦性处理，就从各几何纤维的光滑性判据推出总态射对整个$G$光滑。报告中的双数参数空间反例明确展示此差别。本版显式标出译校补入的步骤：先在$G_{\mathrm{red}}$上工作，再取相关有限型族共同平坦的非空开集$U$，在该开集上构造正确的非光滑坏集并作维数估计。$G$与其约化结构的泛点剩余域相同，所以这一处理保留所研究的泛截面；它不证明族对原来的非约化$G$处处光滑。

\lead{研究计划与已经证明的结果分开。}
作者对Bertini--Zariski推广仍有保留，并给收信人提出继续公理化的建议；本版保留这种未完成的语气。法稿的7.3没有撇号，关于不是8.7(c)直接推论的结果所含否定词则确实保留在原文。上述辨读不填补8.7(c)的历史缺续，也不使所有推广计划成为已经证明的一般定理。新恢复的(2.12)至(2.18)范围句仍是作者的写作概括，不能由本版单向的泛截面光滑性补证宣布其中全部条目已经逐条证明。新增充分条件、反例及译者补证均明确标识，原作者不承担本版后来补入的文字。

\audit{V23-INT03}{继承材料、复现检查和下一入口}{保留原文件并记录实际技术验证的范围}
本版逐字节保留第二十二版四个源码目录中的202个TeX文件，其中包括此前194个继承文件及第二十二版新增8个文件。旧第15节的22条\texttt{V-H-L048}至\texttt{V-H-L069}校注及三条普通脚注仍在完整旧源中，可以核对；新输入链只排印两份回校后的第15节，不重复排印旧入口。旧\path{v_dot_late.tex}的49,224字节、其中第14节34,610字节与第15节14,614字节的精确分拆记录继续保留。Dot69--108的40个模块仍各输入一次。旧报告按版次完整收入后部；其中“本轮”仅指其制作当时。

工程继续使用真实\texttt{ctexbook}的默认字体、字号、版心与双面分页。模块版和独立单文件分别由XeLaTeX构建，检查最终日志，并逐页比较文字、页面尺寸和72dpi像素。新增及受影响页面须实际渲染查看，记录绑定最终PDF与实际图像。具体页数、看图范围和检查结果由随工程的实测文件给出；全书编译通过或抽样版面检查不能代替所有历史数学证明的重新审定。

下一回校点为第二篇第16节，法文物理95页、英文物理62页，已存在的中文入口为\path{text/v1_2.tex}开头的“函数的奇异零点与超奇异零点”。该短篇重编号为V.1（原IV.16），按现有内容次序排在第一篇之前；下一轮应回校其现有译文，而非在第15节后另排一份重复文本。\texttt{EGAV1-2.ps}及其转换稿也是需互校的载体。正式第IV章主体仍从其第一分册物理221页提要与后续章首接续，属于另一未完成分支。

全套书籍尚未完结。第14节既有充分条件、原8.7(c)、8.8、9.5的缺文及13.4(d)的一般基底限制仍保留。历史中7个早期原ZIP容器不可恢复的记录也未改变；保存成员文件不能冒充找回原容器。远端发布须另核实际提交、标签、工作流和附件，不由预设文件名推断成功。
